% WALRUS simulator methods template
% Required packages: amsmath, amssymb, booktabs, tabularx
% The wording is intentionally explicit and can be shortened for the main paper.

\subsection{On-the-fly spectral simulation}
\label{sec:walrus_simulator}

Training spectra were generated on the fly as complex-valued frequency-domain
signals.  For every simulated sample, the network input was decomposed as
\begin{equation}
    x(f) = m(f) + w(f) + \ell(f) + n(f),
    \label{eq:simulator_input}
\end{equation}
where $m$, $w$, $\ell$, and $n$ denote the metabolite, residual-water, lipid,
and receiver-noise components, respectively.  The regression target contained
only the nuisance components,
\begin{equation}
    y(f) = w(f) + \ell(f).
    \label{eq:simulator_target}
\end{equation}
Consequently, subtraction of the predicted nuisance signal from the input
provides the nuisance-suppressed spectrum.  All components were generated or
sampled as complex signals; their real and imaginary parts formed the two input
and output channels of the neural network.

\subsubsection{Metabolite model}

The metabolite signal was synthesized directly from the same processed
LCModel basis functions that were used for spectral fitting.  Let $b_k(t)$ be
the complex FID of basis component $k$ and $c_k\geq0$ its independently sampled
amplitude.  Before line-shape and phase distortions, the metabolite FID was
\begin{equation}
    m_0(t) = \sum_{k=1}^{K} c_k b_k(t).
    \label{eq:metabolite_linear_combination}
\end{equation}
The 7-T basis was prepared at a bandwidth of $2778\,\mathrm{Hz}$, with
$N=840$ complex samples and a proton Larmor frequency of
$297.222931\,\mathrm{MHz}$ ($297.222931\,\mathrm{Hz/ppm}$).  The processed
\texttt{clean\_fid} basis representation was center-cropped in the frequency
domain to the acquisition bandwidth and transformed to the time domain.  Of
the 39 available basis components, 16 were enabled for simulation (Table~
\ref{tab:calibrated_metabolites}).

The original LCModel basis components additionally contained the artificial
reference resonance near $0$~ppm that is used by LCModel during fitting.  This
reference signal is not part of the physiological metabolite spectrum and was
therefore removed before the basis was used for simulation.  Separately for
each basis component, a rank-5 HLSVD model was fitted to the first 8192 samples
of the native FID.  Among the estimated exponentially damped complex
sinusoids, only components with resonance frequencies between $-0.2$ and
$0.2$~ppm were selected.  These components were reconstructed over the full
native FID and subtracted,
\begin{equation}
    b_k^{\mathrm{clean}}(t)
    = b_k^{\mathrm{LCModel}}(t)
    - \sum_{q\in\mathcal{R}_k}
      A_{kq}\exp\!\left(
        \frac{t}{T_{kq}}
        + i\left[2\pi f_{kq}t+\varphi_{kq}\right]
      \right),
    \qquad
    \mathcal{R}_k=\{q:\delta_{kq}\in[-0.2,0.2]\ \mathrm{ppm}\}.
    \label{eq:basis_reference_removal}
\end{equation}
Only $b_k^{\mathrm{clean}}$ was subsequently used in
Eq.~\eqref{eq:metabolite_linear_combination}; the reference-containing basis
itself remained available for the conventional LCModel fitting procedure.

For each sample, a common zero-order phase $\phi_0$, frequency offset
$\Delta f$, Gaussian FWHM $G$, and Lorentzian FWHM $L$ were applied to the
complete metabolite mixture:
\begin{align}
    m_1(t)
    &= m_0(t)
       \exp\!\left[i\left(\phi_0+2\pi\Delta f t\right)\right]
       \exp\!\left[-\pi L t
       -\frac{(\pi Gt)^2}{4\ln 2}\right].
    \label{eq:metabolite_fid_effects}
\end{align}
A first-order phase term was subsequently applied in the frequency domain,
\begin{equation}
    m(f)=\mathcal{F}\{m_1(t)\}\exp(i\phi_1 f),
    \label{eq:first_order_phase}
\end{equation}
where $\phi_1$ is expressed in $\mathrm{rad/Hz}$ and $f=0$ denotes the
spectral center.  The sign conventions in Eqs.~\eqref{eq:metabolite_fid_effects}
and \eqref{eq:first_order_phase} are the conventions used by the simulator.

The total Voigt FWHM $V$ was sampled first.  A line-shape fraction
$r\sim\mathcal{U}(0,1)$ was then used to span the pure-Gaussian to
pure-Lorentzian range.  Specifically,
\begin{align}
    L_{\max} &= \frac{V}{0.5346+\sqrt{0.2166}}, &
    L &= rL_{\max},\\
    G &= \sqrt{(V-0.5346L)^2-0.2166L^2},
    \label{eq:voigt_split}
\end{align}
which satisfies the standard approximation
\begin{equation}
    V = 0.5346L + \sqrt{0.2166L^2+G^2}.
\end{equation}

\subsubsection{Empirical water and lipid components}

Residual-water and lipid line shapes were sampled from subject-specific
in-vivo resource pools rather than generated analytically.  Water was isolated
from the original MRSI FIDs by a 32-component Hankel singular-value
decomposition (HSVD).  Components within the configured water-frequency
window were reconstructed; at 7~T the frequency-offset bounds were
$[-1.35,0.51]$~ppm relative to the water reference.  Isolated-water FIDs from
brain-mask voxels formed the water pool.  Lipid resources were obtained from
voxels in the extracranial lipid mask after subtracting the corresponding
HSVD-derived water FID.

Subjects were sampled uniformly, followed by uniform sampling of resources
within the selected subject.  Water and lipids were drawn from the same subject
for a given synthetic spectrum.  One water spectrum $\widetilde w(f)$ and
$J=10$ lipid spectra $\widetilde\ell_j(f)$ were selected.  The lipid spectra
were combined using positive random weights,
\begin{align}
    u_j &\sim \mathcal{U}(0,1), &
    \alpha_j &= \frac{u_j}{\sum_{q=1}^{J}u_q}, &
    \widetilde\ell(f)&=\sum_{j=1}^{J}\alpha_j
    \widetilde\ell_j(f).
    \label{eq:lipid_mixture}
\end{align}

The global metabolite frequency offset $\Delta f$ was also applied to the
water component because the stored water resources were B0-centered:
\begin{equation}
    \widetilde w_{\Delta f}(f)
    =\mathcal{F}\!\left\{
      \mathcal{F}^{-1}\{\widetilde w(f)\}
      \exp(i2\pi\Delta f t)
    \right\}.
\end{equation}
In contrast, lipid spectra retained their measured frequency distribution and
were not globally shifted.  The simulated metabolite phase terms $\phi_0$ and
$\phi_1$ were not applied to either empirical component; water and lipids
retained their measured complex phase.

Let
\begin{equation}
    A_m=\max_f|m(f)|,\quad
    A_w=\max_f|\widetilde w_{\Delta f}(f)|,\quad
    A_\ell=\max_f|\widetilde\ell(f)|.
\end{equation}
No ppm restriction was used for these maxima.  Sampled dimensionless amplitude
ratios $r_w$ and $r_\ell$ were then imposed as
\begin{align}
    w(f)&=r_w A_m\frac{\widetilde w_{\Delta f}(f)}{A_w},\\
    \ell(f)&=r_\ell A_m\frac{\widetilde\ell(f)}{A_\ell}.
    \label{eq:water_lipid_scaling}
\end{align}
Thus, $r_w$ and $r_\ell$ specify the maximum water-to-metabolite and
lipid-to-metabolite amplitude ratios, respectively.

\subsubsection{Receiver noise and SNR}

Complex white receiver noise was generated with independent unit-variance
Gaussian real and imaginary components.  After applying the acquisition-length
operation described below, its amplitude was adjusted to reproduce the sampled
LCModel-style SNR,
\begin{equation}
    \mathrm{SNR}
    =\frac{\max_f|m(f)|}
    {2\,\mathrm{RMS}\{\Re[n(f)]\}}.
    \label{eq:snr_definition}
\end{equation}
Accordingly, if $n_0$ is an unscaled noise realization, the applied noise scale
was
\begin{equation}
    a_n=\frac{\max_f|m(f)|}
    {2\,\mathrm{SNR}\,\mathrm{RMS}\{\Re[n_0(f)]\}},
    \qquad n(f)=a_n n_0(f).
\end{equation}
This scaling was performed after truncation/zero filling, ensuring that the
requested SNR was realized in the final frequency-domain training spectrum.

\subsubsection{Acquisition length and normalization}

The same acquisition-length operator $\mathcal{A}$ was applied jointly to the
clean full signal, nuisance target, metabolite-only signal, and initially
unscaled noise.  It consisted of an inverse Fourier transform, retention of the
first $N_{\mathrm{acq}}$ FID samples, optional zero filling, and transformation
back to the frequency domain.  For the 7-T experiments reported here,
$N_{\mathrm{acq}}=840$ was fixed and therefore no effective truncation or zero
filling occurred.

The final unnormalized pair was
\begin{align}
    x_{\mathrm{raw}} &= \mathcal{A}\{m+w+\ell\}+n,\\
    y_{\mathrm{raw}} &= \mathcal{A}\{w+\ell\}.
\end{align}
Both were normalized using one scale calculated exclusively from the full
network input,
\begin{align}
    s &= \max_f|x_{\mathrm{raw}}(f)|,\\
    x &= x_{\mathrm{raw}}/s, & y &= y_{\mathrm{raw}}/s.
    \label{eq:shared_normalization}
\end{align}
The identical normalization preserves the relative amplitudes of input and
target.  For the Y-Net ablation, the optional subject-specific L2 lipid
projection was applied to $x_{\mathrm{raw}}$ and normalized by the same $s$ to
form the second input.  Given a matrix $R$ containing the subject's lipid-mask
spectra as rows, the frequency-domain projection operator was
\begin{align}
    M &= R^{\mathrm H}R,\\
    P_{\ell} &= I-(I+\beta M)^{-1},\\
    x_{\mathrm{L2}} &= x_{\mathrm{raw}}P_{\ell}/s.
    \label{eq:l2_projection}
\end{align}
The regularization parameter $\beta$ was selected by bisection such that the
mean absolute diagonal of $(I+\beta M)^{-1}$ equalled $0.99$ within a tolerance
of $10^{-4}$ (maximum 60 iterations).  This second input was used only by the
Y-Net + L2-prior ablation and did not alter the training target.

\subsubsection{Probability distributions}

We used the following reusable distribution families.  A positive calibrated
mixture was defined as
\begin{equation}
    Z_+ \sim \frac{1}{2}\,\mathcal{N}(\mu,\sigma^2)\big|_{Z>a}
    +\frac{1}{2}\,\mathrm{LogNormal}(\mu_{\log},\sigma_{\log}^2)
      \big|_{Z>a},
    \label{eq:positive_mixture}
\end{equation}
where truncation was implemented by inverse-CDF sampling.  A signed calibrated
mixture was
\begin{equation}
    Z_{\pm}\sim
    \begin{cases}
      \mathcal{N}(c,\sigma^2),&p=0.50,\\
      c+\mathrm{LogNormal}(\mu_{\log},\sigma_{\log}^2),&p=0.25,\\
      c-\mathrm{LogNormal}(\mu_{\log},\sigma_{\log}^2),&p=0.25.
    \end{cases}
    \label{eq:symmetric_mixture}
\end{equation}
Here, lognormal parameters refer to the natural logarithm.  A log-uniform
variable $Z\sim\mathrm{LogUniform}(a,b)$ was generated by sampling
$\log Z\sim\mathcal{U}(\log a,\log b)$.

\subsubsection{Calibration of the sampling distributions}

The simulator distributions were calibrated from voxelwise estimates in nine
healthy 7-T MRSI data sets (four Brisbane and five London subjects).  Unless
stated otherwise, all finite values within the respective brain masks were
pooled across subjects.  Thus, the calibration was voxel weighted rather than
subject weighted; no CRLB threshold was applied.  Positive quantities were
additionally restricted to values greater than zero.

Distribution parameters were obtained by robust summary statistics rather
than by maximum-likelihood fitting of a mixture model.  For positive samples
$x_j>0$, the parameters of Eq.~\eqref{eq:positive_mixture} were estimated as
\begin{align}
    \mu &= \operatorname{median}(x_j), &
    \sigma &= \kappa\,\frac{\operatorname{IQR}(x_j)}{1.349},\\
    \mu_{\log} &= \operatorname{median}(\log x_j), &
    \sigma_{\log} &= \kappa_{\log},
        \frac{\operatorname{IQR}(\log x_j)}{1.349}.
\end{align}
For signed samples, the center of Eq.~\eqref{eq:symmetric_mixture} was
$c=\operatorname{median}(x_j)$.  The central Gaussian width was estimated from
$x_j$ as above, whereas the common lognormal distribution for both tails was
estimated from $d_j=|x_j-c|>0$:
\begin{align}
    \sigma &= \kappa\,\frac{\operatorname{IQR}(x_j)}{1.349},\\
    \mu_{\log} &= \operatorname{median}(\log d_j), &
    \sigma_{\log} &= \kappa_{\log},
        \frac{\operatorname{IQR}(\log d_j)}{1.349}.
\end{align}
The left and right tails consequently used identical lognormal parameters and
differed only in sign.  Mixture weights were fixed to $0.5/0.5$ for positive
quantities and $0.5/0.25/0.25$ for the central, positive-tail, and
negative-tail components of signed quantities; they were not fitted to the
data.  The factors $\kappa=\kappa_{\log}=2$ deliberately broadened the robust
empirical widths for training.  In the final 7-T configuration, FWHM was the
exception: its widths were reduced from the initially broadened values
$0.0904$ and $1.4$ to $0.045$ and $0.7$, respectively.

FWHM, SNR, zero-order phase, and first-order phase were calibrated from the
corresponding LCModel maps.  The frequency-shift distribution was calibrated
from the LCModel shift maps of the data before B0 correction, because this is
the frequency variability that the simulator was intended to reproduce.

For metabolite $i$, the LCModel coefficient $c_{i,j}$ in voxel $j$ was first
made relative to the peak magnitude of the complete fitted metabolite signal.
Using the prepared LCModel basis FIDs $b_i(t)$ and the voxelwise LCModel
linewidth $L_j$, this was
\begin{align}
    M_j(\nu) &= \mathcal{F}\!\left\{
        \left[\sum_i c_{i,j}b_i(t)\right]
        \exp(-\pi L_j t)
    \right\},\\
    r_{i,j} &= \frac{c_{i,j}}{\max_{\nu}|M_j(\nu)|}.
\end{align}
The maximum was evaluated over the full spectral axis and, importantly,
$M_j$ contained all fitted metabolites rather than only metabolite $i$.
Positive finite $r_{i,j}$ values were then pooled and calibrated with the same
positive-mixture procedure.

Water and lipid scale ratios were estimated from nine healthy subjects.  The
water FIDs were reconstructed from the unsuppressed-water data.  In the
original calibration used for the trained models, a previously trained
WALINET produced the metabolite-dominated estimate $S_{\mathrm{met}}$ from the
full signal $S_{\mathrm{full}}$.  The nuisance estimate and lipid component
were consequently defined as
\begin{equation}
    S_{\mathrm{nuis}}=S_{\mathrm{full}}-S_{\mathrm{met}},\qquad
    S_{\mathrm{lip}}=S_{\mathrm{nuis}}-S_{\mathrm{water}}.
\end{equation}
The voxelwise ratios were
\begin{equation}
    r_w=\frac{\max_{\nu}|S_{\mathrm{water}}(\nu)|}
              {\max_{\nu}|S_{\mathrm{met}}(\nu)|},\qquad
    r_\ell=\frac{\max_{\nu<4\,\mathrm{ppm}}|S_{\mathrm{lip}}(\nu)|}
                 {\max_{\nu}|S_{\mathrm{met}}(\nu)|}.
\end{equation}
The water and metabolite maxima used the full spectral axis, whereas the lipid
maximum was restricted to frequencies below $4$ ppm.  Finite ratios with a
positive metabolite denominator were pooled and calibrated with the positive
mixture above.  A later L2-based decomposition was used as a WALINET-independent
validation analysis, but it did not define the parameters of the trained
calibrated models.

Table~\ref{tab:calibrated_global_parameters} lists the global 7-T distributions
used for calibrated training.  Parameters specified in ppm were internally
converted to Hz via $1\,\mathrm{ppm}=297.222931\,\mathrm{Hz}$; degrees were
converted to radians, and $\mathrm{deg/ppm}$ to $\mathrm{rad/Hz}$.

\begin{table}[t]
\centering
\caption{Global parameters of the in-vivo-calibrated 7-T simulator.  PM and SM
denote the positive and signed mixtures in Eqs.~\eqref{eq:positive_mixture} and
\eqref{eq:symmetric_mixture}, respectively.}
\label{tab:calibrated_global_parameters}
\begin{tabular}{llll}
\toprule
Parameter & Distribution & Parameters & Unit \\
\midrule
$\Delta f$ & SM & $c=0$, $\sigma=0.116$, $\mu_{\log}=-3.10$, $\sigma_{\log}=2.45$ & ppm \\
$V$ & PM & $\mu=0.056$, $\sigma=0.045$, $\mu_{\log}=-2.88$, $\sigma_{\log}=0.70$, $a=0$ & ppm \\
$\phi_0$ & SM & $c=-12$, $\sigma=50.4$, $\mu_{\log}=2.89$, $\sigma_{\log}=2.10$ & degree \\
$\phi_1$ & SM & $c=1.0$, $\sigma=26.7$, $\mu_{\log}=2.20$, $\sigma_{\log}=2.34$ & degree/ppm \\
SNR & PM & $\mu=9$, $\sigma=5.19$, $\mu_{\log}=2.20$, $\sigma_{\log}=0.5732$, $a=1$ & -- \\
$r_w$ & PM & $\mu=32.4$, $\sigma=36.8$, $\mu_{\log}=3.48$, $\sigma_{\log}=1.10$, $a=0$ & -- \\
$r_\ell$ & PM & $\mu=2.36$, $\sigma=6.65$, $\mu_{\log}=0.857$, $\sigma_{\log}=2.21$, $a=0$ & -- \\
\bottomrule
\end{tabular}
\end{table}

Metabolite amplitudes were sampled independently.  For calibrated training,
each enabled metabolite used Eq.~\eqref{eq:positive_mixture} with $a=0$ and the
parameters in Table~\ref{tab:calibrated_metabolites}.  These amplitudes are
relative basis weights rather than absolute concentrations.

\begin{table}[t]
\centering
\caption{Calibrated 7-T metabolite-amplitude distributions.}
\label{tab:calibrated_metabolites}
\begin{tabular}{lrrrr@{\qquad}lrrrr}
\toprule
Met. & $\mu$ & $\sigma$ & $\mu_{\log}$ & $\sigma_{\log}$ &
Met. & $\mu$ & $\sigma$ & $\mu_{\log}$ & $\sigma_{\log}$ \\
\midrule
Asp  & .381 & .563 & -.964 & 1.48 & Cr   & .525 & .971 & -.644 & 1.58 \\
GABA & .115 & .239 & -2.16 & 2.09 & Glc  & .101 & .244 & -2.29 & 2.29 \\
Gln  & .401 & 1.14 & -.914 & 2.33 & Glu  & 1.11 & 1.55 & .102 & 1.30 \\
GPC  & .220 & .282 & -1.51 & 1.19 & GSH  & .217 & .479 & -1.53 & 1.87 \\
Ins  & .732 & .968 & -.312 & 1.19 & NAA  & 1.38 & 1.07 & .319 & .744 \\
NAAG & .263 & .464 & -1.34 & 1.69 & PCh  & .0953 & .266 & -2.35 & 2.47 \\
PCr  & .539 & .903 & -.619 & 1.50 & Scy  & .0303 & .0452 & -3.50 & 1.40 \\
Tau  & .131 & .328 & -2.03 & 2.19 & 2HG  & .613 & 2.75 & -.490 & 3.38 \\
\bottomrule
\end{tabular}
\end{table}

\paragraph{Heuristic-training ablation.}
The heuristic-training U-Net used the identical signal-formation model but
replaced selected calibrated distributions by broad WALINET-style priors:
\begin{align}
    c_k &\sim \mathcal{N}(1,5^2)\big|_{c_k>0},\\
    \Delta f &\sim\mathcal{U}(-0.5,0.5)\ \mathrm{ppm},\\
    V &\sim\mathcal{U}(0.0135,0.168)\ \mathrm{ppm},\\
    \phi_0 &\sim\mathcal{U}(0,2\pi),\\
    r_w &\sim\mathrm{LogUniform}(0.01,100),\\
    r_\ell &\sim\mathrm{LogUniform}(0.1,1000).
\end{align}
The first-order phase and SNR distributions were unchanged from
Table~\ref{tab:calibrated_global_parameters}.  This ablation therefore isolates
the effect of data-driven simulator calibration while retaining the same basis,
empirical nuisance resources, signal assembly, SNR definition, target, and
normalization.

% Suggested citation placeholders:
% - LCModel basis / LCModel fitting reference
% - HSVD reference
% - Voigt FWHM approximation reference
% - Original WALINET reference for the heuristic ranges
